\chapter{Always, somewhere, someone is fighting for you (Puella Magi Oriko Magica)}
Like Kazumi Magica, the comic series Puella Magi Oriko Magica (by Mura Kuroe, based on a story by Magica Quartet, a collective name for Madoka Magica producer Atsuhiro Iwakami, director Akiyuki Shinbo, writer Gen Urobuchi, and character designer Ume Aoki) must grapple with the trap created by the three inescapable truths of the series—to recap, that witches are both a challenge to the dominant power structure and a corrupt element, that said dominant power structure, represented by Kyubey, is massively unfair, and that said power structure cannot be overthrown within a spinoff, because that overthrow occurs within the final episode of Madoka Magica. Kazumi Magica emphasized the first tenet, and so fell into the trap of vilifying opposition to the system. Oriko Magica instead emphasizes the third tenet, and so avoids that trap, but at the price of becoming a tragedy.

Much like Madoka Magica, Oriko Magica is divided into three arcs, tied together by the titular Oriko. The first is a red herring both in and out of the story: in what is revealed in the third arc to be a ploy to distract Kyubey, Oriko manipulates an abused, orphaned small child named Yuma into contracting as a magical girl. For the reader, this creates the expectation that the comic will largely be—as the first two chapters actually are—the story of the friendship and mentor-pupil relationship between Kyoko and Yuma. However, the second arc belies this expectation, as that story is abruptly cut off at the end of the second chapter; the third and fourth chapters instead follow Mami, as she investigates the serial murders of magical girls and engages in combat with the culprit, Oriko's friend and fellow magical girl Kirika. Despite Kirika's power to slow time, Mami is eventually able to severely injure her, and Kirika flees back to Oriko. The final arc picks up at Mitakihara Middle School, where Oriko and Kirika take over and create a witch's barrier over the entire school. This arc mostly follows Homura, as we learn Oriko's true goal is to save the world from the witch Madoka will eventually become by killing her now, and Homura of course soon confronts and battles Oriko. Kirika becomes a witch (unstated, but implied, is that she began transforming after her battle with Mami, and that the barrier around the school and familiars attacking the students are hers), and Mami, Kyoko, and Yuma fight her. Ultimately, Homura succeeds in killing Oriko, but Oriko kills Madoka in her final moments, causing Homura to reset the timeline.

Tragedy abounds in this story; nearly every character is undone by what appear to be positive impulses. Most obviously, Homura fails to protect Madoka and Oriko dies because they are both determined to protect what is important to them. Oriko cannot risk the destruction of the world by Madoka and thus sees no possible course of action except to try to kill her; Homura cannot risk harm coming to Madoka, and thus sees no possible course of action but destroying Oriko. The end result is that both fail: Homura resets the timeline, ultimately leading to the series timeline, Madoka's rewrite of history, and the destruction of all the timelines where witches existed, including Oriko's; Oriko, meanwhile, kills the Madoka of her timeline, making this yet another failure for Homura. But tragedy abounds for the other characters, as well. Madoka is only in position for Oriko to kill her because of her desire to help Homura; if she'd stayed within the protective barrier Homura made for her, she'd have survived. Kyoko tries to protect Yuma from the dangers of the magical girl life, but in the process makes Yuma worried that she's useless; Yuma then becomes a magical girl in order to be useful to Kyoko. And Kirika wished only to have the courage to stop being a loner, but in the process became a predator.

But perhaps most tragic in all of this is Oriko, who can see the future. The comic ends with a rather ambiguous scene: Kirika approaches a depressed-seeming Oriko, who reveals she has killed a large number of people. Kirika offers to share her burden, and the two go off hand in hand. The Puella Magi Wiki interprets this as another timeline or afterlife \cite{ref105}, but given that the comic has been interspersed throughout with flashbacks to how Oriko and Kirika met and became magical girls, what seems most likely is that this is another such flashback—or, more accurately, the "present" events of the comic are a vision Oriko has just had of her future. She sees that she will kill many, and that she and Kirika will themselves die, but also that this means her timeline's Earth will be spared destruction by Madoka's witch form, and she decides that it's worth it.

If this is what is happening in the final scene, then Oriko's calm determination throughout the series is explained: she is simply acting out her part in a play that has already been written. She knows when and how she will die, and has already decided to embrace her fate for the good of the world. In that sense, she becomes a fallen Madoka, much as Kirika is a fallen Homura. The latter pair share in common that they are utterly devoted to their loved one, willing to do anything for her, but where Homura usually avoids hurting humans or magical girls and instead focuses on fighting witches and Kyubey, Kirika hunts magical girls and only fights witches to directly protect Oriko. Similarly, both Oriko and Madoka are willing to sacrifice themselves to save the world, but Oriko is willing to kill, while Madoka's goal is to save everyone.

This, then, is the Oriko Magica solution to the trap: the unfair system of witches and magical girls cannot be broken within the comic, so focus on a smaller effort to save the world. Show how everyone is doomed by the system around them, and depict their failure to overthrow the system as a tragedy for all involved. However, the comic evades a second trap that might have opened here, that of implying that the system is unbreakable, by deliberately showing Homura resetting and ending on the relatively happy note of Kirika and Oriko sharing their burdens.

Homura continues to work toward the ending we see in the series. In the meantime, comfort can be found in one another. Though all magical girls will someday be broken by Kyubey's system, so long as it stands, as Yuma says, "Someday is not today."
